\section{Related work}\label{sec:related}

\subsection{Quality-aware multi-video summarization}

Extractive summarization selects existing frames or temporal segments, whereas
abstractive approaches generate derived representations
~\citep{truong2007video_abstraction,apostolidis2021survey}. Multi-video methods
add explicit cross-source redundancy and diversity. Video-MMR formalized the
setting~\citep{li2010videommr}; sparse and collaborative formulations improved
diversity across topic-related videos~\citep{panda2017diversity,
panda2017collaborative}; graph models subsequently represented inter-video
relations directly~\citep{wu2020mvsgcn}. These methods address
representativeness, coverage and redundancy, but they do not validate a
visual-quality proxy used to choose between redundant candidates.

HSMVideoSumm instead derived selection criteria from human summarization
practice and implemented each criterion as an explicit pipeline stage
~\citep{barbieri2021content,barbieri2021tese}. This modularity permits a
component-level audit: the same clusters, filters and ordering can be held fixed
while the operation used for visual quality is tested and replaced. The object
of comparison is therefore neither a new summarizer nor an overall benchmark
against SumMe or TVSum; it is one inspectable decision inside a multi-source
news pipeline.

The use of image-quality assessment in summarization is not itself new.
VSQUAL used objective IQA measures for shot-boundary detection, keyframe
extraction and summary construction~\citep{cirne2014vsqual}. Jeong et al. also
combined representativeness with contrast, blur and image-skew measures when
selecting keyframes for static video summaries~\citep{jeong2016static}. The present
contribution is different: it does not claim novelty for adding IQA, but for a
mechanistic construct-validation audit of an already deployed quality operation,
followed by dependency-aware propagation of an operation-level substitution.

\subsection{No-reference quality assessment and transfer}

Image quality assessment predicts perceived degradation. Full-reference
measures require an undistorted original and are unsuitable when candidates come
from different broadcasters~\citep{wang2004ssim,zhai2020iqa_survey}. NR-IQA
instead estimates quality from the test image alone. BRISQUE and NIQE model
departures from natural-scene statistics~\citep{mittal2012brisque,
mittal2013niqe}; DBCNN and HyperIQA learn opinion-supervised convolutional
representations~\citep{zhang2020dbcnn,su2020hyperiqa}; MUSIQ and TOPIQ use
multi-scale or top-down attention~\citep{ke2021musiq,chen2024topiq}.

Vision--language formulations based on Contrastive Language--Image Pre-training
(CLIP) provide a different inductive bias. CLIP-IQA uses
antonym prompts~\citep{wang2023clipiqa}; LIQE jointly predicts quality,
distortion and scene~\citep{zhang2023liqe}; QualiCLIP learns from rankings of
synthetically degraded pairs without opinion scores~\citep{agnolucci2024qualiclip}.
Q-Align maps multimodal-model outputs to discrete quality levels
~\citep{wu2024qalign}. These models differ not only architecturally but also in
their exposure to opinion databases. A result on training data measures fit,
whereas use in a new broadcast domain requires transfer. Cross-database
evaluation is therefore essential, and published work has shown that
in-database NR-IQA performance can overstate generalization
~\citep{pinson2023nrbroadcast,cao2024datacentric}.

KADID-10k supplies controlled synthetic distortions
~\citep{lin2019kadid}; KonIQ-10k and LIVE In the Wild (CLIVE) contain authentic
in-the-wild images~\citep{hosu2020koniq,ghadiyaram2016clive}. Their differences
make them complementary validation settings, but comparisons among unmatched
databases and checkpoints do not identify causal effects of training exposure
or distortion regime. The complete descriptive contrasts are consequently kept
in Supplementary Table~S15 rather than presented as a principal contribution.

Two further cautions affect the broadcast application. First, professional and
user-generated video can contain temporal impairments that a frame-wise image
score cannot observe; video-quality models were developed partly for this reason
~\citep{wang2021uvq,huang2023xgcvqa}. Second, high correlation on a reused
database need not imply difficult pairwise discrimination. Database range and
subjective-score uncertainty affect correlation~\citep{krasula2016methodology},
and recent multimodal benchmarks report that comparative quality judgments can
behave differently from absolute scoring~\citep{zhang2024qbenchplus}. Newer
distributional quality models also exceed earlier reference-model results on
some databases~\citep{you2025deqascore}. Accordingly, Q-Align and DOVER are
treated here as complementary contemporary references, not definitive ground
truth or evidence that temporal modeling itself causes an advantage.

\subsection{The unresolved measurement problem}

SIFT, BoVW and TF-IDF encode local-feature occurrence and rarity for retrieval
~\citep{lowe2004sift,sivic2003videogoogle,csurka2004bovw,
salton1988tfidf}. Spatial Information similarly quantifies scene complexity
rather than perceptual quality~\citep{itu_p910,robitza2021siti}. Such features
may correlate with opinion scores in a particular collection, but correlation
does not show that they measure quality or remain valid under controlled
degradation.

The gap addressed here is narrower than the general problem of quality-aware
summarization. Existing work establishes multi-video selection criteria,
specialized NR-IQA methods and prior use of IQA in summarization, yet the
operation implementing HSMVideoSumm's human-grounded visual-quality criterion
was not validated as a quality measure. This study links a mechanistic analysis
of that operation to external human scores, construction-defined degradation
pairs and a controlled operation-level substitution in the pipeline. That
combination distinguishes construct validation from merely integrating an
available IQA model.
